\documentclass[letterpaper,twocolumn]{article}

\usepackage[margin=0.75in]{geometry}
\usepackage{times}
\usepackage{helvet}
\usepackage{courier}
\usepackage{amsmath}
\usepackage{amssymb}
\usepackage{booktabs}
\usepackage{graphicx}
\usepackage{multirow}
\usepackage{url}
\usepackage[hidelinks]{hyperref}
\usepackage{algorithm}
\usepackage{algorithmic}
\usepackage{microtype}
\usepackage[numbers,sort&compress]{natbib}

\title{Safe-Teacher Gated Projection-Aware PPO for Actuator-Feasible Full-Car Active Suspension Control}

\author{Anonymous Authors}

\date{}

\begin{document}

\maketitle

\begin{abstract}
Deep reinforcement learning (DRL) is attractive for active suspension control because it can adapt to nonlinear road--vehicle interactions without requiring a perfectly identified plant model.
Yet direct continuous-action RL is poorly matched to safety-critical suspension deployment: random online exploration can create excessive body acceleration, unsafe suspension travel, actuator-rate violations, and a mismatch between sampled policy actions and the force actually executed after actuator and safety projection.
This paper proposes Safe-Teacher Gated Projection-Aware PPO (STG-PPO), an offline-to-online control framework for full-car active suspension.
The controller first initializes PPO by behavior cloning a safe passive teacher, then constrains online PPO through three deployment-aware mechanisms: a delta-parameterized action space tied to the previously executed force, projection-consistency and actuator-feasibility regularization, and a state-dependent safe-teacher improvement gate that permits active force only when road preview excitation, body attitude acceleration, and safety margin indicate that deviation from the passive teacher is justified.
On a MuJoCo full-car benchmark with physically calibrated vehicle parameters, correlated ISO-style roads, pothole-roll, sine-roll, and mixed-road scenarios, STG-PPO improves over PPO from scratch on all core safety and feasibility metrics.
In the full 200-episode matrix, STG-PPO reduces unsafe steps from 106.0 to 0.0, action-delta RMS from 142.92 N to 13.98 N, actuator tracking RMS from 225.09 N to 20.86 N, and policy projection error from 0.3626 to 0.0001.
Against SAC, STG-PPO obtains higher return and lower body/pitch acceleration and projection error, while SAC remains closer to passive behavior in roll and actuator activity.
These results support a focused claim: safe-teacher gated online PPO is a viable and auditable path to actuator-feasible active-suspension adaptation, rather than an unconstrained replacement for classical or off-policy controllers.
\end{abstract}

\section{Introduction}

Active suspension systems can improve ride comfort and handling by applying controlled vertical forces between the vehicle body and wheels.
Modern perception and road-preview systems further make proactive suspension control possible, but they also introduce a difficult learning-control problem: the controller must exploit preview information, respect suspension travel and tire-load limits, and remain feasible under actuator lag and force-rate constraints.
Model predictive control (MPC) with road preview is a natural solution because it can encode constraints explicitly \citep{papadimitrakis2022preview_mpc,li2024speed_preview_mpc}.
However, accurate full-car models, online optimization cost, perception delay, and unmodeled actuator dynamics make purely model-based control brittle in deployment.

Deep RL provides an appealing alternative because continuous-action actor-critic methods such as PPO \citep{schulman2017ppo}, TD3 \citep{fujimoto2018td3}, and SAC \citep{haarnoja2018sac} can learn nonlinear feedback policies in simulation.
Recent work has begun to apply DRL and demonstrations to active suspension \citep{tan2024active_demo}.
The obstacle is not simply sample efficiency.
For active suspension, the dangerous regime is early online adaptation: an untrained policy can issue large actuator forces, rely on hidden safety projection, or produce high-frequency commands that cannot be tracked by the physical actuator.
In this setting, a high return after long training is not sufficient; the learning process itself must be conservative, auditable, and aligned with the actually executed control input.

This paper develops Safe-Teacher Gated Projection-Aware PPO (STG-PPO), a conservative online PPO framework for full-car active suspension.
The central idea is that PPO should not directly own the actuator command.
Instead, it should begin near a safe teacher, operate in a feasible incremental action space, and earn deviations from passive suspension only when the current state indicates a plausible improvement opportunity.
Concretely, STG-PPO combines offline behavior cloning, delta-parameterized actions, safety projection, projection-aware PPO penalties, and a safe-teacher improvement gate.
The gate is not merely a reward term: it changes the deployed policy class by blending PPO commands toward the passive teacher whenever road excitation or safety margin is insufficient.

The contributions are:
\begin{itemize}
    \item We formulate active suspension online adaptation as a projection-aware PPO problem in which the learning objective penalizes mismatch between sampled policy actions and actuator-feasible executed actions.
    \item We introduce a safe-teacher improvement gate that converts PPO from an always-active force generator into a state-dependent residual over a passive teacher.
    \item We provide a full-car MuJoCo benchmark with actuator lag, rate limits, road preview uncertainty hooks, safety metrics, and decomposed reward logging.
    \item We report ablations showing that STG-PPO repairs unsafe PPO exploration and is competitive with TD3/SAC under physically calibrated full-car dynamics, while preserving actuator feasibility.
\end{itemize}

\section{Related Work}

\paragraph{Continuous-control RL.}
DDPG introduced deep deterministic actor-critic learning for continuous action spaces \citep{lillicrap2015ddpg}.
TD3 reduces actor-critic overestimation through clipped double critics, delayed actor updates, and target smoothing \citep{fujimoto2018td3}.
SAC optimizes a maximum-entropy objective and is often robust and sample-efficient in continuous control \citep{haarnoja2018sac}.
PPO is an on-policy method with clipped policy-ratio updates \citep{schulman2017ppo}; its conservative update is attractive for online adaptation, but the standard formulation does not know whether sampled actions will later be modified by actuator or safety projection.
STG-PPO addresses this missing deployment layer.

\paragraph{Learning from demonstrations.}
Demonstration-augmented RL can reduce unsafe or inefficient early exploration.
Deep Q-learning from demonstrations showed that expert data can accelerate RL in domains where learning from scratch is costly \citep{hester2018dqfd}.
Active suspension work has also explored expert demonstrations, for example by pretraining DRL from PID samples \citep{tan2024active_demo}.
Our setting differs in two ways: the teacher is not treated as a final policy to imitate forever, and deviation from the teacher is controlled online by a state-dependent improvement gate.

\paragraph{Active suspension and road preview.}
Preview MPC for active suspension explicitly optimizes future comfort and handling under road-preview information \citep{papadimitrakis2022preview_mpc,li2024speed_preview_mpc}.
This line is strong when the model and optimizer are reliable, but computational cost and model mismatch remain concerns.
Recent DRL approaches consider nonlinear uncertainty, expert demonstrations, and actuator constraints in active suspension \citep{tan2024active_demo}.
STG-PPO is complementary: it keeps the adaptive capacity of DRL while explicitly constraining how the learned policy may depart from a safe teacher under actuator limits.

\section{Problem Formulation}

We consider a full-car vertical dynamics environment with sprung-body heave, pitch, and roll; four unsprung wheel masses; passive suspension and tire elements; and four active force commands.
At time $t$, the agent observes
\begin{equation}
    o_t = [x_t, p_t, u^{\mathrm{exec}}_{t-1}],
\end{equation}
where $x_t$ includes vehicle states, $p_t$ is road preview, and $u^{\mathrm{exec}}_{t-1}$ is the previously executed actuator force.
The policy samples a normalized action $a_t \in [-1,1]^4$.
The environment executes a force after actuator and safety processing:
\begin{equation}
    u^{\mathrm{exec}}_t = \Pi_{\mathcal{U}(x_t)}(u^{\mathrm{cmd}}_t),
\end{equation}
where $\Pi_{\mathcal{U}(x_t)}$ enforces force magnitude, rate feasibility, and safety-margin constraints.

The reward is decomposed as
\begin{equation}
r_t = -(J_{\mathrm{comfort}} + J_{\mathrm{handling}} + J_{\mathrm{act}} + J_{\mathrm{safety}} + J_{\mathrm{dev}}),
\end{equation}
with comfort penalties on body heave, pitch, and roll accelerations; handling penalties on suspension and tire-load terms; actuator penalties on force magnitude, force delta, command delta, tracking error, saturation, and energy proxy; and optional deviation from a teacher or prior.
We evaluate policies using return, unsafe steps, body/pitch/roll acceleration RMS, action-delta RMS, actuator tracking RMS, saturation ratio, and projection error.

\section{Method}

\subsection{Offline Teacher Initialization}

Before online PPO, we collect safe teacher trajectories from a passive controller.
The actor is behavior-cloned with
\begin{equation}
    \mathcal{L}_{\mathrm{BC}}(\theta)
    =
    \mathbb{E}_{(o,u^E)}
    \left[
    \left\|
    \pi_\theta(o) - u^E / u_{\max}
    \right\|_2^2
    \right],
\end{equation}
where $u^E=0$ for the passive teacher in the current implementation.
This gives the online policy a safe initial behavior and makes subsequent deviations measurable.

\subsection{Delta-Parameterized Projection-Aware Action Space}

Instead of interpreting PPO output as an absolute actuator force, STG-PPO interprets it as a bounded increment from the previous executed force:
\begin{equation}
    \tilde{u}_t
    =
    \mathrm{clip}\left(
        u^{\mathrm{exec}}_{t-1}
        +
        \eta u_{\max} \tanh(a_t),
        -\rho u_{\max}, \rho u_{\max}
    \right),
\end{equation}
where $\eta$ controls the maximum per-step force increment and $\rho$ bounds the action magnitude.
The projected execution is
\begin{equation}
    u^{\mathrm{exec}}_t=\Pi_{\mathcal{U}(x_t)}(\tilde{u}_t).
\end{equation}
We log the normalized projection error
\begin{equation}
    e^\Pi_t =
    \frac{\|\tilde{u}_t-u^{\mathrm{exec}}_t\|_2}{u_{\max}\sqrt{d_u}},
\end{equation}
which measures how much hidden correction is required before the policy becomes feasible.

\subsection{Safe-Teacher Improvement Gate}

The key algorithmic addition is a state-dependent gate $g_t \in [0,1]$ that blends the PPO action toward the passive teacher:
\begin{equation}
    u^{\mathrm{cmd}}_t
    =
    g_t \tilde{u}_t + (1-g_t)u^{E}_t.
\end{equation}
In the current passive-teacher setting, $u^E_t=0$.
The gate is computed from road preview excitation, body attitude acceleration demand, and safety margin:
\begin{align}
    g^{p}_t &= \mathrm{clip}\left(
        \frac{\mathrm{RMS}(p_t)-\tau_p}{\sigma_p},0,1
    \right),\\
    g^{a}_t &= \mathrm{clip}\left(
        \frac{A_t-\tau_a}{\sigma_a},0,1
    \right),\\
    g_t &= \mathrm{clip}\left(
        g_{\min} + (g_{\max}-g_{\min})
        \max(g^{p}_t,g^{a}_t)g^{s}_t,
        g_{\min},g_{\max}
    \right),
\end{align}
where $A_t$ combines heave, pitch, and roll acceleration magnitudes, and $g^s_t$ shrinks the gate near safety-limit margins.
Thus PPO acts aggressively only when the state contains evidence that active force is useful and safe.

\subsection{Projection-Aware PPO Objective}

PPO optimizes the clipped surrogate objective
\begin{equation}
    \mathcal{L}_{\mathrm{clip}}
    =
    \mathbb{E}_t
    \left[
    \min(\rho_t A_t,\mathrm{clip}(\rho_t,1-\epsilon,1+\epsilon)A_t)
    \right],
\end{equation}
where $\rho_t$ is the policy likelihood ratio.
We add feasibility penalties:
\begin{equation}
    \mathcal{L}
    =
    \mathcal{L}_{\mathrm{clip}}
    -
    \lambda_\Pi \mathbb{E}_t[\rho_t e^\Pi_t]
    -
    \lambda_U \mathbb{E}_t[\rho_t c^{\mathrm{unsafe}}_t]
    -
    \lambda_\Delta \mathbb{E}_t[\rho_t \Delta u_t].
\end{equation}
The penalties discourage actions that become feasible only because the safety layer silently changes them.

\begin{algorithm}[t]
\caption{STG-PPO Online Adaptation}
\label{alg:stgppo}
\begin{algorithmic}[1]
\STATE Collect safe teacher trajectories $\mathcal{D}_E$
\STATE Initialize actor by minimizing $\mathcal{L}_{\mathrm{BC}}$
\FOR{each online episode}
    \STATE Reset environment and obtain $o_0$
    \FOR{each step $t$}
        \STATE Sample PPO action $a_t \sim \pi_\theta(\cdot|o_t)$
        \STATE Convert $a_t$ to delta force $\tilde{u}_t$
        \STATE Compute safe-teacher gate $g_t$
        \STATE Blend command $u^{\mathrm{cmd}}_t=g_t\tilde{u}_t+(1-g_t)u^E_t$
        \STATE Execute $u^{\mathrm{exec}}_t=\Pi_{\mathcal{U}(x_t)}(u^{\mathrm{cmd}}_t)$
        \STATE Log reward components, projection error, unsafe flag, and gate
    \ENDFOR
    \STATE Update PPO using projection-aware objective
\ENDFOR
\end{algorithmic}
\end{algorithm}

\section{Experimental Setup}

\paragraph{Environment.}
Experiments use a MuJoCo full-car active-suspension environment \citep{todorov2012mujoco}.
The model includes heave, pitch, roll, four unsprung masses, four road anchors, passive spring-damper elements, and active suspension force actuators.
The calibrated benchmark uses a 1500 kg sprung body, 2160 kg m$^2$ pitch inertia, 460 kg m$^2$ roll inertia, 59 kg unsprung mass per wheel, front/rear suspension stiffness of 35000/38000 N/m, and tire stiffness of 190000 N/m.
The benchmark includes class-B and class-C correlated random roads, mixed full-road excitation, pothole-roll, and sine-roll scenarios.
Road preview is available to the controller, while clean preview is retained for diagnostics and ablation.

\paragraph{Compared methods.}
We compare: (i) PPO from scratch with raw-force actions and the same downstream safety projection; (ii) BC-PPO with offline imitation, delta action parameterization, projection-aware penalties, and the safe-teacher gate; (iii) residual BC-PPO variants around a full-car MPC-lite prior; (iv) TD3; (v) SAC; and (vi) passive and MPC-lite reference controllers.
For fair off-policy baselines, TD3 and SAC do not inherit PPO-specific action parameterization or safe-teacher gating.

\paragraph{Metrics.}
Higher return is better.
For all other core metrics, lower is better: unsafe steps, body/pitch/roll acceleration RMS, action-delta RMS, actuator tracking RMS, saturation ratio, and policy projection error.
We also report the mean safe-teacher gate to measure how often PPO departs from passive behavior.

\section{Results}

\subsection{Repeated-Seed PPO Ablation}

Table~\ref{tab:seed} reports BC-PPO minus PPO-scratch deltas across three seeds.
STG-PPO supports the core claim on every seed: return improves, unsafe steps decrease, all three acceleration RMS metrics improve, action-delta and tracking errors decrease, and projection error drops.

\begin{table}[t]
\centering
\caption{Repeated-seed evidence. Return delta is higher-is-better; other deltas are lower-is-better.}
\label{tab:seed}
\resizebox{\columnwidth}{!}{
\begin{tabular}{lrrrrrrrr}
\toprule
Seed & Return & Unsafe & Body & Pitch & Roll & $\Delta u$ & Track & Proj. \\
\midrule
42 & 6326.75 & -20.80 & -1.85 & -2.31 & -3.84 & -158.11 & -235.92 & -0.4800 \\
43 & 11580.64 & -52.40 & -2.81 & -3.16 & -6.41 & -151.98 & -226.77 & -0.4406 \\
44 & 10346.37 & -112.40 & -2.32 & -1.04 & -2.99 & -90.57 & -135.15 & -0.3260 \\
\bottomrule
\end{tabular}}
\end{table}

\subsection{Full-Matrix Comparison}

Table~\ref{tab:full} summarizes learned-controller means for the physically calibrated 200-episode matrix.
Compared with PPO from scratch, STG-PPO reduces unsafe steps from 106.0 to 0.0 and projection error from 0.3626 to 0.0001.
The mean gate is 0.1878, indicating that the policy remains near passive on easy segments and opens mainly under stronger excitation.

\begin{table*}[t]
\centering
\caption{Full-matrix learned-controller means. Gate is the mean safe-teacher improvement gate.}
\label{tab:full}
\begin{tabular}{llrrrrrrrrr}
\toprule
Variant & Ctrl. & Return & Unsafe & Body & Pitch & Roll & $\Delta u$ & Track & Proj. & Gate \\
\midrule
PPO scratch & PPO & -8806.61 & 106.00 & 1.033 & 6.546 & 2.617 & 142.92 & 225.09 & 0.3626 & 1.000 \\
STG-PPO & PPO & -423.91 & 0.00 & 0.480 & 0.748 & 0.080 & 13.98 & 20.86 & 0.0001 & 0.188 \\
Residual BC-PPO & PPO & -9838.13 & 65.60 & 2.151 & 1.209 & 9.488 & 162.69 & 245.60 & 0.3166 & 1.000 \\
Safe residual BC-PPO & PPO & -13598.16 & 105.40 & 0.613 & 0.947 & 11.920 & 163.97 & 246.40 & 0.3349 & 1.000 \\
TD3 & TD3 & -523.07 & 0.00 & 0.626 & 0.887 & 0.626 & 122.10 & 182.18 & 0.0421 & 1.000 \\
SAC & SAC & -577.53 & 0.00 & 0.499 & 0.751 & 0.054 & 0.00 & 0.00 & 0.1880 & 1.000 \\
\bottomrule
\end{tabular}
\end{table*}

\paragraph{Against SAC and TD3.}
STG-PPO is not claimed to dominate all off-policy methods.
In the calibrated e200 run, STG-PPO has higher mean return than both TD3 and SAC and lower body/pitch acceleration and projection error.
However, SAC behaves close to a passive low-action policy: it has lower roll acceleration and essentially zero actuator activity, while STG-PPO opens the gate on a small fraction of states.
We therefore interpret the off-policy comparison as evidence of a tradeoff: STG-PPO is the stronger projection-aware online-adaptation policy, while SAC and passive control remain strong low-action references on the current road set.

\subsection{Ablation Interpretation}

The residual-prior branches remain weak in the current implementation.
This negative result is useful: a physics prior is not automatically beneficial if the prior controller is itself poorly matched to the full-car benchmark or if the residual action space is not constrained as carefully as the direct delta PPO branch.
The paper therefore centers on the direct safe-teacher gated PPO mechanism and treats residual MPC/LQR priors as future work.

\section{Discussion}

The results support a narrow but important claim.
For active suspension, online PPO from scratch is not acceptable because it can issue unsafe or untrackable actuator commands.
Adding only scalar reward penalties is insufficient when the policy likelihood is computed on actions that differ from the executed projected force.
STG-PPO changes both the learning objective and deployed policy class: the policy is initialized near a safe teacher, parameterized as an incremental force, penalized for projection mismatch, and gated by state-dependent evidence for improvement.

The current gate is hand-designed and interpretable.
This is useful for debugging and deployment audits, but a stronger CCF-A contribution would replace it with a learned improvement certificate.
For example, a passive-vs-active advantage critic could estimate a lower confidence bound on the benefit of opening the gate under road-preview uncertainty.
The gate would then be learned from data while preserving the safety principle: do not deviate from the safe teacher unless improvement is likely.

\section{Limitations}

First, the evidence is simulation-only and still covers a limited set of road distributions and random seeds.
Second, passive control remains the best nominal reference on the calibrated road set, so the current claim is actuator feasibility and online PPO repair, not active-control superiority over passive suspension.
Third, SAC remains competitive as a low-action baseline; harder road classes and payload/parameter variation are needed to show when active force is essential.
Fourth, the current gate uses hand-tuned thresholds.
Fifth, the related-work section should be expanded with additional verified recent active-suspension DRL papers before submission.

\section{Conclusion}

We presented STG-PPO, an offline-imitation-initialized, projection-aware, safe-teacher-gated PPO framework for full-car active suspension control.
The method addresses a deployment-specific problem in continuous-action RL: the policy should not be trained as if raw sampled actions are executed when actuator and safety projection substantially modify them.
Across the calibrated full-time matrix, STG-PPO improves over PPO from scratch on return, safety, comfort, action smoothness, actuator tracking, and projection error.
It also outperforms TD3/SAC on return and projection consistency, while passive and SAC remain strong low-action references.
The next step is to learn the improvement gate as an uncertainty-aware certificate and to evaluate on harder road distributions where passive control is insufficient.

\bibliographystyle{plainnat}
\bibliography{references}

\end{document}
